\documentclass[openany, oneside]{book}

\newcommand{\lib}[1]{\input{../../Libs/LaTeX/#1}}

\lib{Text}
\lib{Math}
\lib{Russian}

\newcommand{\vf}{\vspace{5mm}}

\usepackage{tempora}

\usepackage{epigraph}

\usepackage{geometry}
\geometry{a4paper, top=30mm, left=30mm, right=30mm, bottom=35mm}

\renewcommand{\epigraphsize}{\normalsize}
\renewcommand{\footnotesize}{\normalsize}

% \renewcommand{\epigraphsize}{\large}
% \renewcommand{\footnotesize}{\large}

% \hyphenpenalty = 10000
% \fontdimen2\font=0.7ex
% \fontdimen3\font=0.4em
\renewcommand{\baselinestretch}{1.015}

\sloppy
\begin{document}

\thispagestyle{empty}

\emptysymbol

\vspace{70mm}

\includegraphics{pictures/small_world.png}

\vspace{10mm}

\ \ \ \ \ \ {\LARGE \bfseries КАЛЕНДАРЬ СНОВИДЕНИЙ}
% \ \ \ \ \ \ {\LARGE \bfseries Жизнь и сновидения — страницы одной и той же книги.}

\vspace{7mm}

\ \ \ \ \ \ {\LARGE Роман Максимович}

\large

\newpage

\emptysymbol

\vspace{2cm}

\begin{flushleft}
{\it
<<I can almost feel you walking in the distance,

But I know that it's not reality\dots

If you were echoing around inside these parallel dimensions all along,

Know I'll find you in my dreams\dots>>
}

\vf

--- Lost Sky, <<Dreams, pt. II>>

\vf
\vf

\rule{\textwidth}{.7pt}

\vf
\vf

{\it <<Ты пришел из мира Ямы

Радостно ты, мудрый, используешь смертных

\dots

Вначале тебя видело всепоглощающее дно

\dots

Оттуда, о сон, ты пришел сюда,

Глубоко пряча свою форму>>}

\vf

--- <<Атхарваведа>>, учение о Благородном Восьмеричном Пути Будды.

\vf
\vf

\rule{\textwidth}{.7pt}

\vf
\vf

\textbf{Предисловие}\\


Это рассказ о прошлом и будущем мира, в котором люди способны воплощать сновидения в действительность. Во избежание разрушительных инцидентов места проживания людей покрываются <<нейринами>> --- онейронными ингибиторами, которые полностью блокируют сновидения. Однако же, всякое правило имеет исключения\dots \

Всё в мире связано, ибо мир тесен.

Перед прочтением рекомендуется ознакомиться с моим предыдущим текстом, <<Sonnik>>, первым в серии <<Календарь сновидений>>.

\end{flushleft}

\fontdimen2\font=0.35em
\fontdimen3\font=0.1em

\chapter{Познание в практике}

Они штурмовали завод по сборке онейронных ингибиторов. Их было четверо, объединённых стремлением дать миру вновь увидеть сны. Тише самой ночи, крепкие и гибкие, одетые в маскировочные костюмы, они быстро прокрадывались через заросли ко входу в сооружение завода, бесшумно снимая охранников из винтовок\dots

Что, поверили? Нет, огранизованная преступная группировка сновидцев <<Грёзы Бхагавана>> всем составом пребывала в маленькой съёмной квартирке в пригороде Лос Анджелеса. Хью развалился в кресле у камина, мирно похрапывая, Зак и Шерон в обнимку спали на диване, привалившись друг к другу. Завершением этой уютной картины был лидер отряда, Алан Ривер. Он лежал на спине прямо на полу, закинув ноги наверх, на диван. По его словам, так кровь приливала к мозгу, и ему лучше думалось. На стене ровно мигала диодами коробочка активного <<нейрина>>.\\

Крепкие и гибкие, одетые в маскировочные костюмы, они стояли на зелёном холме, а перед ними высилась громада производственного комплекса. Небо было окрашено закатными ядовито-зелёными тонами, местами оно замутнялось бордовыми кучевыми, рубиновыми перистыми и алыми слоистыми облаками. До самого горизонта рисовался густой лесной массив глубоководных оттенков, с частыми лазурными островками.

--- Yes! Наконец-то получилось!

--- Два месяца, мать их за ногу!

--- Чёрт, как же я заколебался рыть инфу по этому треклятому заводу, пока моё упрямое подсознание не согласится его \textit{воплотить}!

Хью, сорокалетний мужик, прыгал от счастья, пока Зак и Шерон не заключили его в радостные объятия. Алан стоял с гордой улыбкой на лице. Он поднял руку в небо, привлекая внимание команды. Громкий щелчок пальцами, и в воздухе заиграла легендарная \href{https://www.youtube.com/watch?v=btPJPFnesV4}{<<Eye of the Tiger>>}.

--- Друзья! Свершим же то, зачем мы с вами здесь собрались!

Зак еле сдержал смешок.

--- Шеф, и где ты только такого лексикона набрался? <<Свершим же>>\dots \ Но ты прав, --- парень обернулся к зданию завода, --- погнали!\\

--- Боб, ты это видел?

--- Ты о чём? Чёрт, я не выспался, что ли\dots \ Слушай, подежуришь за меня чуть-чуть? Совсем в сон клонит\dots

Оба свалились на землю, как подкошенные.\\

--- Зак! На тебе охрана. Всё как планировали.

--- Понял!

Медленно, словно неохотно, от завода к ближайшему посёлку потянулся косяк спящих людей, парящих в воздухе, но обвисших, как мокрые тряпки. Среди них были охранники, чиновники и прочий персонал комплекса. <<Грёзам Бхагавана>> ни к чему были жертвы. Километр за километром, люди удалялись от здания, а потом ложились кучкой неподалёку от населённого пункта.

Шерон, Хью и Алан разбежались в разные части производственного сооружения. Каждый член отряда направлялся к одному из трёх термоядерных энергетических установок, питающих работу комплекса. Добежав до реактора, Шерон принялась судорожно вспоминать, как её учил Хью. В попытке устранить недоверие граждан к <<нейринам>> американское государство сделало их производство предельно прозрачным. Каждый шаг, каждая технология, вся информация о процессе изготовления ингибиторов была в открытом или почти открытом доступе. Поэтому Хью, обладая достаточными познаниями в ядерной физике, вполне мог спланировать взрыв энергетической установки завода. Сначала Шерон вскрыла панель управления и отключила аварийную систему защиты реактора. Далее она извлекла из активной зоны управляющие графитовые стержни, а потом погрузила реактор в надкритическое состояние. Проследив рост реактивности на сенсорах, Шерон удовлетворённо кивнула и ринулась к ближайшему выходу. Алан и Хью проделали ту же процедуру, пока Зак заканчивал эвакуацию, и уже через минуту все четверо стояли на лазурно-голубом холме. Они переглянулись, после чего Алан с размаху вмазал Заку прямо по лицу, и парень исчез, не успев издать ни звука. Та же судьба послигла и Хью. Алан посмотрел на Шерон и виновато покачал головой: прости, мол, не могу. Шерон манерно закатила глаза, подошла к пятиметровому обрыву, которым заканчивался холм, а затем с весёлым криком спрыгнула вниз.

Один из главных вопросов, которые задаёт себе человек в осознанном сновидении, заключается в том, как из этого сновидения, собственно, выйти. Как проснуться? Опытный сновидец знает множество способов. Самый простой и доступный из них --- болевой шок. Нужно, однако, быть осторожным, чтобы полученный во сне ушиб или перелом вдруг не \textit{воплотился}. В конце концов, умрёшь во сне --- умрёшь и в реальности. Тело не живёт без мозга. Алану, в отличие от его товарищей, такой экстрим был не по душе. Он закрыл глаза и ловко сел в <<позу лотоса>>, скрестив ноги.

--- Намо Ариясвапная Буддхая... --- сосредоточенно провозгласил он, и было ясно, что он вкладывает глубокий смысл, все свои мысли в каждое слово. Миг, и лазурно-голубой холм оказался пуст.

А потом одновременно прогремели три взрыва.\\

На следующий день по телевизору шёл репортаж.

<<Вчера, 20 июня 2114 года, в районе 6 часов вечера был совершён теракт в штате Миссури, в районе города Сент-Луис. В результате преступных действий был полностью уничтожен завод по производству онейронных ингибиторов, также известных как <<нейрины>>. Однако, к удивлению органов расследования, теракт не унёс ни одной человеческой жизни. Весь персонал комплекса был обнаружен в ближайшем посёлке местными жителями. Полиция сообщает, что преступники не оставили никаких видимых следов своего присутствия. Эксперты, в свою очередь, предполагают, что нападение имело онейронный характер, но мнения о том, как именно злоумышленники осуществили теракт, всё ещё расходятся. С вами была Кейт Локуэлл, следите за подробностями инцидента после рекламной паузы.>>

\begin{center}
    * * *
\end{center}

--- <<Так я слышал. Однажды Бхагаван обитал в Шравасти, в роще Джета, в саду Анатхапиндады (Дающий Защиту Обездоленным), поддерживая и распространяя Дхарму. А пребывал он в своём пребывании вместе с великим собранием бхикшу, числом в двенадцать с половиной сотен, с прославленными старейшинами, великими учениками, всеми архатами. Там были старейшины Шарипутра, Махамаудгальяна, Махакашьяпа, Махакаппхина, Махакатьяяна, Махакауштхила, Шуддхипинтхака, Гавампатин, Бхарадваджа, Калодаийн, Ваккула и Анирудха; и вместе с этими другие великие ученики, сопровождаемые бодхисаттвами-махасаттвами. Там были бодхисаттвы Манджушри, Гандхавастин, Авалокитешвара (Созерцающий Звуки Мира), Нитьёдукта и Аникшиптадхура; и вместе с этими другие бодхисаттвы-махасаттвы. Также там были Шакра, царь богов, и Брахма, господин этого мира Саха; и вместе с ними сотни тысяч миллионов божественных существ.

В это время Бхагаван обратился к старейшине Шарипутре: есть, Шарипутра, на востоке, через десятки тысяч миллионов буддхакшетр, одна буддхакшетра, называемая Страной Благородных Сновидений. В ней как раз сейчас пребывает Будда Татхагата по имени Ариясвапна, архат, полностью самопросветлённый, утвердившись в тонкой и чудесной Дхарме. Как ты думаешь, Шарипутра, по какой причине эта земля называется Страной Благородных Сновидений? В этой стране нет у живых существ ни телесных, ни душевных страданий, ибо пребывают они в состоянии полного духовного отстранения. Как ты думаешь, Шарипутра, что позволяет живым существам достичь полного духовного отстранения? Дхарму вспомнив, они погружаются в сон. Будду вспомнив, они погружаются в сон. Сангху вспомнив, они погружаются в сон. И каждое из десятков тысяч миллионов живых существ внутри себя свою буддхакшетру создаёт. По этой причине, Шарипутра, эта земля и называется Страной Благородных Сновидений.

Далее, Шарипутра, Страна Благородных Сновидений окружена семью рядами ограждений, семью рядами деревьев, семью рядами сетей с колокольчиками, и все они, прекрасные и красивые, сделаны из драгоценностей четырёх видов; а именно из золота, серебра, ляпис-лазури и хрусталя; такими прекрасными качествами и достоинствами, Шарипутра, украшена эта буддхакшетра.

Далее, Шарипутра, Страна Благородных Сновидений окружена шестью Нефритовыми Вратами с шести сторон; а именно с севера, востока, юга, запада, с верхней стороны и с нижней стороны. И лишь только обладая сильным духом, может любое живое существо отворить одно из этих Нефритовых Врат, осознанно или неосознанно, для себя или для других живых существ; такими прекрасными качествами и достоинствами, Шарипутра, украшена эта буддхакшетра.

Далее, Шарипутра, в Стране Благородных Сновидений небо сделано из ядовито-зелёных изумрудов, а облака сделаны из алых рубинов, а земля сделана из глубоководных сапфиров, прекрасных и красивых. Далее, Шарипутра, там, когда дуют приятные освежающие ветры, ряды драгоценных деревьев и драгоценных сетей с колокольчиками прекрасные звуки издают. Такими качествами и достоинствами, Шарипутра, украшена эта буддхакшетра.

Далее, помнишь ли ты, Шарипутра, что в Стране Благородных Сновидений каждое живое существо в состоянии полного духовного отстранения пребывает, и каждое внутри себя свою буддхакшетру создаёт? Видишь ли, Шарипутра, ежели в одной из этих буддхакшетр некое событие или шевеление Дхармы произойдёт, оно тут же в Стране Благородных Сновидений \textit{воплотится} и за её пределами распространится. Например, ежели одно живое существо в своём духовном отстранении цветение лотоса наблюдает, в Стране Благородных Сновидений тут же цветок лотоса вырастает, наполняя Дхармой, наполняя Буддой, наполняя Сангхой всех живых существ. Такими прекрасными качествами и достоинствами, Шарипутра, украшена эта буддхакшетра.

Далее, Шарипутра, в этой земле Татхагаты Ариясвапны живут чистые бодхисаттвы; те, кто уже не вернётся и связанные одним рождением. И, Шарипутра, число тех бодхисаттв нелегко измерить, только в неизмеримых, неисчислимых множествах оно может быть посчитано.

Как ты думаешь, Шарипутра, по какой причине этот Татхагата называется Ариясвапна (Благородный Сон)? Именно потому, Шарипутра, что этот Татхагата беспрестанно пребывает во сне, и сон его поистине благороден. По этой причине этот Татхагата и называется Ариясвапна. И, Шапипутра, прошло десять кальп с тех пор, как этот Татхагата достиг наивысшего и полного просветления и стал Буддой.

Так же как сейчас я, Шарипутра, имя Татхагаты Ариясвапны восхвалял, так и на востоке Татхагата Акшобхья, Татхагата Подобный Меру, Татхагата Великая Меру, Татхагата Сияния Меру, Татхагата Прекрасное Знамя, все Будды Бхагаваны на востоке, бесчисленные, как песок в Ганге, каждый в своей собственной буддхакшетре, сие учение распространяют и провозглашают.

Так же в южной стороне света Татхагата Лунный и Солнечный Светильник, Татхагата Прекрасный Свет, Татхагата Великих Светящихся Плеч, Татхагата Светильник Меру, Татхагата Бесконечная Доблесть, все Будды Бхагаваны на юге, бесчисленные, как песок в Ганге, каждый в своей собственной буддхакшетре, сие учение распространяют и провозглашают.

Так же в западной стороне света Татхагата Неизмеримая Жизнь, Татхагата Бесконечных Плеч, Татхагата Бесконечных Свойств, Татхагата Великий Свет, Татхагата Великий Драгоценный Свет, Татхагата Сияния Чистого Света, все Будды Бхагаваны на западе, бесчисленные, как песок в Ганге, каждый в своей собственной буддхакшетре, сие учение распространяют и провозглашают.

Так же в северной стороне света Татхагата Великих Светящихся Плеч, Татхагата Звучания Музыки Литавров, Татхагата Всеобщего Звучания, Татхагата Солнечнорожденный, Татхагата Неприступный, Татхагата Сияние Водяного Водоёма, Татхагата Солнце, все Будды Бхагаваны на севере, бесчисленные, как песок в Ганге, каждый в своей собственной буддхакшетре, сие учение распространяют и провозглашают.

Так же в нижней стороне света Татхагата Лев, Татхагата Прекрасный Свет, Татхагата Дхарма, Татхагата Охраняющий Дхарму, Татхагата Подобный Дхарме, все Будды Бхагаваны нижнего направления, бесчисленные, как песок в Ганге, каждый в своей собственной буддхакшетре, сие учение распространяют и провозглашают.

Так же в верхней стороне света Татхагата Голос Брахмы, Татхагата Созвездие Царя, Татхагата Царь, Подобный Знамени Индры, Татхагата Наивысшего Аромата, Татхагата Ароматный Свет, Татхагата Великих Светящихся Плеч, Татхагата С Телом, Покрытым Драгоценными Цветами, Татхагата Царь Дерева Сала, Татхагата Красоты Драгоценного Цветка, Татхагата Различающий Все Значения, Татхагата Счастье Горы Меру, все Будды Бхагаваны верхнего направления, бесчисленные, как песок в Ганге, каждый в своей собственной буддхакшетре, сие учение распространяют и провозглашают: вы должны верить в эту сутру, в которой все Будды восхваляются за их непостижимые качества и достоинства, сутру, всеми Буддами принимаемую и защищаемую.

Эти слова Будды Бхагавана услышав, и Шарипутра, и все бхикшу, и все бодхисаттвы вместе со всем миром богов, людей, асуров и гандхарвов, возрадовались и Будду восхвалили.>>

--- Что это? --- недоуменно спросил я. --- Звучит, конечно, очень смешно\dots \ <<Татхагата Бесконечных Плеч>>\dots

--- Это Малая Ариясвапна-сутра буддийской школы Благородных Сновидений, --- спокойно ответил Муген, и голос его был предельно серьёзным, --- я прочитал её тебе с той целью, чтобы ты понял, что \textit{обратное воплощение}, как ты его называешь, существовало во все времена.

Мы с ним сидели под старым, даже древним, но очень красивым сливовым деревом, стоящим посреди широкого луга под лазурным небом. Я --- и мой друг по имени Токигава Муген. Фамилия Токигава пишется иероглифами <<время>> и <<река>>, а имя Муген --- иероглифами <<мечта>> (<<сон>>) и <<реальность>>. Также его можно прочитать как <<бесконечность>>. Муген долго втолковывал мне сложную этимологию своего имени, но чётко я понял лишь одно: значение даётся неспроста, оно обладает глубоким смыслом и пророчит судьбу.\\

Наши встречи происходят весьма спонтанно. Бывает, я не вижу его годами, а бывает, он является в мои сны каждую ночь. На нём обычно широкие штаны, скорее похожие на юбку, два слоя кимоно и красивое хаори, всюду покрытое узором бегущей по горному склону реки. Я же сижу в простой футболке и джинсах.\\

--- Сколько тысячелетий разделяет нас с тобой? --- спросил он, когда мы встретились в первый раз.

--- Ровно одна тысяча лет.

--- Воистину, неподвластна времени Страна Благородных Сновидений!\\



\end{document}